Throughout $n(x) = 1/(e^{x} + 1)$, $u = e^{-y}$, $\theta := r\,d/dr$, and all series are series of analytic functions of $r > 0$ whose terms and termwise derivatives are bounded by $C(1 + m)^{2}e^{-cm}$ locally uniformly in $r$ ($m$ the summation index), so they may be differentiated termwise. The array's Killing flow acts on mode $k$ by $e^{-J\omega_{k}t}$ with $J = e_{2}e_{3}$ the boundary modes' unit; at $r = 1$ one window gives $e^{-2\pi J(k + 1/2)} = -1$, the half turn (the check verifies $J^{2} = -1$). The theta functions are scalar complex analysis with numbers in $\mathrm{span}\{1, I_4\}$.

\emph{The functions.} The modes of $\mathcal{A}(r)$ have $x_{k} = 2\pi(k + \tfrac12)/r$; those with $k \ge 0$ have $x = 2\pi a_{n}/r$, $n = k + 1$, and those with $k < 0$ are their particle-hole images. Above the filled sea each pair contributes $2\ln(1 + e^{-x})$ to $\ln Z$ and $2xn(x)$ to $\beta E$. The zeta-regularised symmetric sea energy at spacing $\kappa/r$ and offset $\tfrac12$ is $\tfrac{\kappa}{2r}B_{2}(\tfrac12) = -\kappa/(24r)$ (Theorem \claimref{synthesis.null-ordering-of-counted-units-is-zeta-symmetric}(ii)), that is $-\pi/(12r)$ in units of $1/\beta$. Hence $f$, $e$ and $e_{\mathrm{sym}}$ as stated, and $S = \ln Z + \beta E = f + e_{\mathrm{sym}}$, which does not depend on the ordering. With $\beta$ varied at fixed $L$, $r = L/\beta$ and $\beta\,\partial_{\beta} = -\theta$, so $\beta E_{\mathrm{sym}} = -\beta\,\partial_{\beta}\ln Z = \theta f$; termwise, $\theta\ln(1 + e^{-x}) = xn(x)$ because $\theta x = -x$, so $\theta f = -\pi/(12r) + e(r) = e_{\mathrm{sym}}(r)$, consistent with the definition. The continuum's energy values: by Theorem \claimref{synthesis.null-ordering-of-counted-units-is-zeta-symmetric}(i) the density is $\kappa^{2}/(48\pi)$ per complex component, so over $L = r\beta = 2\pi r/\kappa$ the energy is $\kappa r/24$ and $\beta$ times it is $\pi r/12$.

\emph{(i)} By clause (a) of Fact \claimref{synthesis.imported-theta-function-identities}, $\vartheta_{3}(\tau)/\eta(\tau) = q^{-1/24}\prod_{n \ge 1}(1 + q^{n - 1/2})^{2}$. At $\tau = I_4/r$, $q = e^{-2\pi/r}$ and the logarithm of the right side is $f(r)$. By clause (b), $\vartheta_{3}$ and $\eta$ acquire the same factor $\sqrt{-I_4\tau}$ under $\tau \mapsto -1/\tau$, so the quotient is invariant, and $-1/(I_4/r) = I_4r$. Hence $f(r) = f(1/r)$, and writing out $f(1/r)$ gives $\pi r/12 + c(r)$. The value at $r = 1$ is the stored $\tfrac12\ln 2$ (Theorem \claimref{synthesis.half-integer-weight-array-carries-half-ln2}(i), whose sum of mode entropies is $f(1) + e_{\mathrm{sym}}(1)$ with $e_{\mathrm{sym}}(1) = 0$ by (ii)). The check confirms the duality at $r = 2, 4, \tfrac12, \tfrac45$ and the value at $r = 1$ to $10^{-30}$ by rational interval arithmetic.

\emph{(ii)} Apply $\theta$ to the dual form: $\theta(\pi r/12) = \pi r/12$ and $\theta\ln(1 + e^{-y}) = -yn(y)$ for $y = 2\pi ra_{n}$, which gives the stated formula. From $f(r) = f(1/r)$, $f'(r) = -r^{-2}f'(1/r)$, so $e_{\mathrm{sym}}(r) = rf'(r) = -(1/r)f'(1/r) = -e_{\mathrm{sym}}(1/r)$, and $e_{\mathrm{sym}}(1) = 0$. Sign: at fixed $L$ the symmetric-ordered energy is $E_{\mathrm{sym}} = re_{\mathrm{sym}}(r)/L$ with $\beta = L/r$, and $dE/d\beta = -\operatorname{Var}(E) < 0$ in a Gibbs state with nonconstant energy, so $re_{\mathrm{sym}}(r)$ is strictly increasing in $r$; it vanishes at $r = 1$, hence the signs, and $r = 1$ is the only zero. Comparison (a) is $1/r = r$. For (b), $D(r) = e_{\mathrm{sym}}(r) + \pi/(12r) - \pi r/12$ by the definition of $e_{\mathrm{sym}}$; inserting the dual formula gives the stated expression, and both $e_{\mathrm{sym}}$ and $\pi/(12r) - \pi r/12$ are odd under $r \mapsto 1/r$, so $D$ is. For $r \ge 1$ write $12rD(r)/\pi = 1 - G(r)$ with $G(r) = 48\sum_{n}a_{n}r^{2}/(e^{y_{n}} + 1)$. Each term has $r$-derivative $a_{n}r\,[2(e^{y} + 1) - ye^{y}]/(e^{y} + 1)^{2}$, negative because $y \ge \pi > 2 + 2e^{-\pi}$. So $G$ is strictly decreasing on $[1, \infty)$, and $G(1) = 1$ because $D(1) = e_{\mathrm{sym}}(1) = 0$. Hence $D > 0$ on $(1, \infty)$, $D < 0$ on $(0, 1)$ by oddness, and $r = 1$ is the only zero. (c) is the statement on $e_{\mathrm{sym}}$. At $r = 1$ the formula for $e_{\mathrm{sym}}$ reads $\pi/12 = 4\pi\sum_{n}a_{n}/(e^{2\pi a_{n}} + 1)$, which is the stated series identity, and it is equivalent to $D(1) = 0$. Fixed spacing: the array of spacing $\kappa$ at inverse temperature $\beta' = \beta/r$ has $x_{k} = \beta'\kappa(k + \tfrac12) = 2\pi(k + \tfrac12)/r$, so $\beta'E_{\mathrm{sym}} = e_{\mathrm{sym}}(r)$. A window $2\pi/\kappa$ is $r$ thermal periods at that temperature, so the continuum's boost-ordered value is $\beta'E = \pi r/12$ by the definition of the extensive value at window length $r$ (the density $\tfrac{\pi}{12}T^{2}$, which at $T = \kappa/2\pi$ is the stored $\kappa^{2}/(48\pi)$ and at general $T$ is the limit $\lim_{R}e(R)/R = \pi/12$ per thermal period, from the dual formula); the passage from the boost ordering to the null ordering subtracts the fixed amount $\kappa/24$ per window (clause (i) of the stored theorem), that is $\pi/(12r)$ in units of $1/\beta'$. The difference is $e_{\mathrm{sym}}(r) - \pi r/12 + \pi/(12r) = D(r)$ in units of $1/\beta' = r\kappa/2\pi$.

\emph{(iii)} $S = f + \theta f = (rf)'$. In the dual form $rf = \pi r^{2}/12 + rc(r)$, so $S = \pi r/6 + c + \theta c$, and $\theta c = -2\sum y_{n}n(y_{n})$: this is $\Delta$. Put $t(y) = \ln(1 + u) - yu/(1 + u)$. From $u - u^{2}/2 \le \ln(1 + u) \le u$ and $u(1 - u) \le u/(1 + u) \le u$,
\[
-(y - 1)u - \tfrac12u^{2} \;\le\; t(y) \;\le\; -(y - 1)u + yu^{2} .
\]
\emph{Sign for $r \ge \tfrac12$.} $t(y) < 0$ whenever $y > 1 + u$ (from $\ln(1 + u) \le u$), and $y_{n} \ge \pi/2 > 1 + e^{-\pi/2}$, so every term of $\Delta$ is negative. \emph{Sign for $r < \tfrac12$.} Use the direct form: $S(r) = 2\sum_{k \ge 0}s(x_{k})$, $s(x) = \ln(1 + e^{-x}) + xn(x)$ the entropy of one mode (Theorem \claimref{synthesis.boundary-mode-thermal-entropy}), with $x_{k} = (2k + 1)x_{0}$, $x_{0} = \pi/r > 2\pi$, and $s(x) \le (1 + x)e^{-x}$. Since $1 + x_{k} \le (2k + 1)(1 + x_{0})$, $S(r) \le 2(1 + x_{0})e^{-x_{0}}\sum_{k}(2k + 1)e^{-2kx_{0}} < \tfrac{201}{100}(1 + x_{0})e^{-x_{0}}$, while $\pi r/6 = \pi^{2}/(6x_{0})$. The function $x(1 + x)e^{-x}$ is decreasing for $x \ge 2$ and at $x = 2\pi$ is below $\tfrac{9}{100}$, so $\tfrac{201}{100}x_{0}(1 + x_{0})e^{-x_{0}} < \tfrac15 < \pi^{2}/6$ and $S(r) < \pi r/6$. \emph{Remainder for $r \ge 1$.} The first term gives $2t(y_{1}) = -2(\pi r - 1)e^{-\pi r} + \rho_{1}$ with $-u^{2} \le \rho_{1} \le 2yu^{2}$, $y = \pi r$, so $|\rho_{1}| \le 2\pi r\,e^{-2\pi r}$. For $n \ge 2$, $|2t(y_{n})| \le 2y_{n}e^{-y_{n}}$ (for $y \ge 1$ both parts of $t$ are at most $yu$), and $\sum_{n \ge 2}2y_{n}e^{-y_{n}} = 2\pi r\,e^{-3\pi r}\sum_{j \ge 0}(2j + 3)e^{-2\pi rj} < \tfrac{61}{10}\pi r\,e^{-3\pi r} < \tfrac{3}{10}\pi r\,e^{-2\pi r}$, using $e^{-\pi} < \tfrac{1}{23}$. Together $|\rho| < 3\pi r\,e^{-2\pi r}$. The value at $r = 1$ is $S(1) - \pi/6$ with $S(1) = \tfrac12\ln 2$. The limit $S(R)/R \to \pi/6$ follows. That it equals the density integral: $S(R)/R = \tfrac{1}{2\pi}\sum_{k}h\,s(h(k + \tfrac12))$ with $h = 2\pi/R$ is a Riemann sum of $\tfrac{1}{2\pi}\int_{\R}s$, and for an integrable function of bounded variation the Riemann sum differs from the integral by at most $h$ times the total variation, here $2\ln 2$ ($s$ is even and decreases from $\ln 2$ to $0$ on $x \ge 0$); so $\int_{\R}s(x)\,dx = \pi^{2}/3$. The interval values are the check's. \emph{The Clausius form.} $S = f + \theta f$ gives $S' = f' + (\theta f)' = \tfrac1r\bigl(e_{\mathrm{sym}} + re_{\mathrm{sym}}'\bigr) = \tfrac1r(re_{\mathrm{sym}})'$, the first law $dS = \beta'\,dE$ at fixed window; and $rD = re_{\mathrm{sym}} - \pi r^{2}/12 + \pi/12$ gives $\tfrac1r(rD)' = \tfrac1r(re_{\mathrm{sym}})' - \pi/6 = \Delta'$. As $r \to 0$, $S(r) \to 0$ (the direct-form bound above), so $\Delta(r) \to 0$ and $\Delta$ is the stated integral.

\emph{(iv)} $\beta^{2}\operatorname{Var}(E) = \beta^{2}\partial_{\beta}^{2}\ln Z = ((\beta\partial_{\beta})^{2} - \beta\partial_{\beta})\ln Z = (\theta^{2} + \theta)f$. Direct form: $(\theta^{2} + \theta)(\pi/(12r)) = 0$ and, with $\theta x = -x$ and $n' = -n(1 - n)$, $(\theta^{2} + \theta)\ln(1 + e^{-x}) = x^{2}n(1 - n)$: the sum of the modes' variances, independent of the ordering. Dual form: $(\theta^{2} + \theta)(\pi r/12) = \pi r/6$ and, with $\theta y = y$, $(\theta^{2} + \theta)\ln(1 + e^{-y}) = y^{2}n(1 - n) - 2yn$. For $y \ge \pi$, $y(1 - n(y)) = y/(1 + e^{-y}) \ge \pi/(1 + e^{-\pi}) > 3$, so each dual term is positive for $r \ge 1$; the dual sum is at most $2\sum y_{n}^{2}e^{-y_{n}}$, which tends to $0$, so $v(R)/R \to \pi/6$, and the Riemann-sum argument of (iii) identifies the limit with the density integral ($x^{2}n(1 - n)$ is integrable and of bounded variation). The value at $r = 1$ is half the stored value per counted unit (Theorem \claimref{synthesis.counted-unit-energy-at-half-offset}(iv)); the check re-bounds it and verifies the dual form at $r = 1, 2$.

\emph{(v)} Theorem \claimref{synthesis.residue-from-mode-entropy}(i) gives $\kappa_{c} = 1/(4\langle s\rangle)$ per counted unit under the area-law hypothesis. A unit of $\bar u$ components carrying $\sigma$ each has $\langle s\rangle = \bar u\sigma$: $\sigma = \tfrac12\ln 2$ gives $1/(2\bar u\ln 2)$ (the stored statement, Theorem \claimref{synthesis.the-residue-requires-two-arrays-per-unit-on-average-and-fixes-no-grouping}(i), (ii)), $\sigma = \pi/6$ gives $3/(2\pi\bar u)$. The ratio $3\ln 2/\pi$ does not involve $\bar u$. Equality with $1/(4\ln 2)$ needs $\bar u = 6\ln 2/\pi$, which the check places in $(1.3238, 1.3239)$. For arrays of window $r$, one array carries $S(r)$ per $r$ modular periods, so $\langle s\rangle$ per component per modular period is $S(r)/r$ and the components per Planck area per modular period are $r/(4S(r))$; by (iii) this tends to $3/(2\pi)$. On the cosmological horizon $|J| = 4\pi\kappa_{c}N^{2}$, and $4\pi \cdot 3/(4\pi) = 3$. The hypothesis is the area-law condition, so, as the cited theorem says of itself, this is a conditional reformulation and not an admissible derivation of the residue.

\emph{Status.} (i) to (iv) are proved given the imported Fact; their numerical values are certified by rational interval arithmetic. (v) is arithmetic under the stated hypothesis. Nothing here assumes Premise (Q), (N$_{+}$) or the window reading, and nothing here decides which entropy the count should use.
